\title{Large Language Models Can Automatically Engineer Features for Few-Shot Tabular Learning}

\begin{abstract}
Large Language Models (LLMs) possess outstanding abilities to address complex and novel reasoning tasks, presenting significant opportunities within tabular learning—a domain critical for various practical scenarios. In this work, we introduce a new in-context learning paradigm called \textsf{FeatLLM}, wherein LLMs are utilized as feature constructors to craft enhanced representations of tabular input data optimized for predictive modeling. The constructed features are then leveraged by a straightforward machine learning algorithm, such as linear regression, to assess class probabilities, thereby enabling highly effective few-shot learning. Crucially, at inference, \textsf{FeatLLM} exclusively employs the discovered features with the simple predictive model, obviating the need to query the LLM for each individual test instance. Additionally, \textsf{FeatLLM} only demands access to the LLM's API—circumventing prompt length restrictions—and avoids repeated LLM interactions during deployment. Empirical evaluations across a broad spectrum of tabular benchmarks from different domains reveal that \textsf{FeatLLM} generates superior rules, leading to notable improvements in predictive performance (with average gains exceeding 10\%) compared to prior approaches such as TabLLM and STUNT.
\end{abstract}

\section{Introduction}
Large Language Models (LLMs) have recently demonstrated remarkable prowess in producing appropriate outputs for tasks not previously encountered during training, without necessitating bespoke fine-tuning for individual tasks~\cite{creswell2022selection,imani2023mathprompter,taylor2022galactica}. Owing to their training on large and diverse text datasets, LLMs encapsulate substantial prior knowledge, allowing them to supplant traditional expert input across numerous domains~\cite{Genomes2021,singhal2023large,wilcox2003role}, and supporting automation in a variety of use cases, such as conversation analysis~\cite{finch2023leveraging} and program synthesis or correction~\cite{fan2023automated}. Notably, when provided with task specifications and a handful of exemplars via prompts, LLMs demonstrate an ability to interpret task context, selecting pertinent features and examples with contextual attention~\cite{brown2020language,zhao2023survey}. Such observations point towards the viability of employing LLMs in the role of expert feature engineers capable of discerning salient characteristics within data.

LLMs’ capabilities to reason and exploit background knowledge are now being harnessed in tabular learning scenarios. Consider, for example, prior facts such as the correlation between advanced age and increased susceptibility to specific diseases, which are valuable in medical prediction settings. Contemporary research has thus taken the approach of articulating tasks and feature explanations in natural language, linearizing tabular entries, and presenting this structured information to LLMs~\cite{dinh2022lift,wang2023anypredict}. Downstream, these models are applied through strategies like in-context learning~\cite{nam2023semi} or efficient fine-tuning methods~\cite{dinh2022lift,hegselmann2023tabllm}. The integration of LLM-provided knowledge has delivered robust improvements, especially for scenarios with limited supervision, consistently surpassing traditional tabular modeling baselines~\cite{hegselmann2023tabllm}.

Existing methods for tabular learning utilizing large language models (LLMs) exhibit notable drawbacks. Typically, each sample requires a separate inference from an LLM for end-to-end prediction, significantly increasing computational demands. Additionally, these approaches often rely on fine-tuning the LLM to achieve competitive accuracy, which is impractical for LLMs lacking publicly available training infrastructure or tuning APIs. This is especially pertinent as many state-of-the-art LLMs only provide restricted access through inference APIs~\cite{anil2023palm,openai2023gpt4}. Moreover, handling lengthy prompts remains a challenge for most LLM-based frameworks~\cite{liu2023lost}. As the dimensionality of tabular data increases and its serialized textual representation becomes longer, performance tends to deteriorate or the methods cease to be applicable. Collectively, these issues substantially hinder the effective integration of LLM-based tabular learning in applied settings.

To address these limitations, our proposed method redefines the utility of LLMs, moving away from direct end-to-end predictions. We instead task LLMs with the generation of new features, capitalizing on their knowledge to augment feature engineering. We present a novel in-context learning methodology, termed \textsf{FeatLLM}, which is designed to capture the implicit "criteria" used by LLMs when making predictions. For example, in a disease classification scenario, the LLM can articulate rules mapping feature conditions to the presence or absence of the disease. These derived criteria, or rules, are then utilized to construct fresh features that can supplant traditional features for downstream tasks. This paradigm shift enables simpler models for subsequent inference, thereby reducing latency once the feature generation is completed. By exploiting the LLMs’ in-context learning capabilities, relevant features can be extracted simply by providing example demonstrations within the prompt, obviating the need for model retraining. Additionally, employing techniques such as feature bagging from ensemble learning~\cite{breiman2001random} helps alleviate prompt length issues, allowing the methodology to scale more gracefully with high-dimensional tabular datasets.

\textsf{FeatLLM} organizes the process of generating novel features through prompt engineering into two distinct stages: (1) comprehending the underlying task and deducing how features relate to target variables; (2) leveraging insights from this analysis, along with selected few-shot examples, to establish definitive rules for class separation. This systematic reasoning mechanism encourages LLMs to disregard irrelevant features when formulating these rules. Once the rules are developed, they are executed on the dataset (as programmatic instructions), converting the samples into binary features that indicate conformity to each rule. Subsequently, a simple, low-complexity model is fitted to the newly constructed features, providing class probability estimates. To bolster reliability, an ensemble methodology is used, involving repeated feature generation and feature bagging. By fine-tuning the number of features and sample subsets included in the bagging step, \textsf{FeatLLM} circumvents the issue of oversized prompts. The method does not require training and incurs minimal inference costs, supporting widespread application of LLMs across a diverse range of real-world tasks.

Our experiments span 13 distinct tabular datasets under limited data scenarios, demonstrating both the effectiveness and consistency of \textsf{FeatLLM}. We show that LLMs skillfully integrate both domain-specific prior knowledge and information gleaned from the data itself to craft high-quality features. Across multiple benchmarks and settings, our framework surpasses existing few-shot learning alternatives. The implementation is accessible via an anonymized GitHub repository at~\texttt{https://github.com/Sungwon-Han/FeatLLM}.

\section{Related Work}

\subsection{Few-Shot Learning with Tabular Data} Recent developments in few-shot learning have made it feasible to extract broadly applicable representations from data, while incurring only limited annotation overhead~\cite{chen2018closer}. Although the methodology was initially pioneered in the context of image analysis, its effectiveness has been established for other modalities such as natural language, speech, and, more recently, structured tables~\cite{majumder2022few,nam2023stunt,schick2021exploiting}. Nevertheless, extracting knowledge from just a handful of labeled points is prone to the problem of overfitting to incidental relationships within the data~\cite{bartlett2021deep}. To address this, contemporary studies have sought to enhance inductive bias during training by tapping into supplementary information sources or by injecting prior, domain-specific insights. Strategies include making use of large-scale but unlabeled tabular datasets paired with appropriate augmentation strategies in a semi-supervised learning framework~\cite{ucar2021subtab,yoon2020vime}, or relying on unsupervised pre-training regimes that do not depend on the final downstream task~\cite{somepalli2021saint}. Common objectives for these unsupervised approaches encompass partial masking or introducing noise to data entries, accompanied by reconstruction tasks~\cite{arik2021tabnet,majmundar2022met}, or augmenting data to facilitate contrastive learning paradigms~\cite{bahri2022scarf,chen2020simple}. Yet, due to the inherent variability and heterogeneity present in tabular data, it remains difficult to define transformations that generalize, and these methods tend to yield limited improvements when applied in few-shot scenarios~\cite{nam2023stunt}.

A growing body of research has turned to model pre-training across a diverse array of tabular datasets—irrespective of their direct relevance to any particular target task—and subsequent transfer learning~\cite{zhang2023generative,levin2022transfer,zhu2023xtab}. A case in point, TransTab~\cite{wang2022transtab}, leverages transformers to learn higher-level semantic dependencies between columns by leveraging column names along with raw table contents. The model thereby amasses transferrable knowledge from disparate datasets, which can be utilized for novel tasks in a zero-shot or few-shot fashion. Alternatively, TabPFN~\cite{hollmann2022tabpfn} employs synthetic data generation by mixing distribution families to train a Bayesian neural network; at inference time, the target task’s train and test samples are provided without additional model updates. Knowledge garnered in this pre-training phase confers a performance edge over standard tree-based baselines, especially under few-shot constraints.

\subsection{Language-Interfaced Tabular Learning} A distinct and promising strategy for integrating prior knowledge leverages the emergent capabilities of large language models (LLMs)~\cite{anil2023palm,openai2023gpt4}. Benefiting from training on expansive, heterogeneous textual resources, LLMs have shown exceptional generalization to previously unseen problems, which has led to their adoption for tasks spanning multiple domains~\cite{brown2020language}. Recent explorations in the context of tables have involved adapting LLMs by encoding the tabular data into a text-based format before passing them into model prompts~\cite{dinh2022lift,hegselmann2023tabllm,wang2023anypredict}. Incorporating not only the data entries but also supplementary meta-information such as feature and task descriptors enables LLMs to model dependencies between predictors and targets for downstream inference. Some studies develop these capabilities further by engaging in parameter-efficient model adjustments through techniques such as LoRA~\cite{hu2021lora} or IA3~\cite{liu2022few}, whereas others employ an in-context learning approach—supplying several demonstration examples within the prompt itself, obviating the need for additional tuning~\cite{dinh2022lift,hegselmann2023tabllm,nam2023semi,slack2023tablet}.

Although previous approaches have shown success in advancing few-shot tabular learning, their reliance on conducting inference entirely via the LLM complicates deployment in industrial settings. To address this, the present work shifts away from the traditional end-to-end black-box paradigm in which every sample's prediction depends on an LLM. Instead, the focus is placed on having the LLM infer explicit class-wise rules or conditions. In \textsf{FeatLLM}, these derived rules are converted into programmatic representations, allowing for the creation of new features. These features facilitate prediction independently of the LLM, utilizing in-context learning to formulate rules based on existing domain knowledge and limited few-shot data.

\section{Methods}
\textsf{FeatLLM} operates by identifying ``rules"—that is, criteria relevant to each answer class—based on the provided few-shot instances, as illustrated in Figure~\ref{fig:main_model}. The LLM extracts these rules, which are subsequently transformed into binary indicators for each sample, denoting rule satisfaction. These binary features then undergo a dot product with a set of trainable, non-negative weights to quantify the probability associated with each class. The training procedure, described in Section~\ref{Sec:training}, enables learning the influence of each feature on class prediction. The methodology incorporates multiple iterations of bagging, followed by ensemble aggregation. Further elaboration on the problem setup and each procedural step is presented below.

\vspace*{-0.12in}
\paragraph{Problem formulation.} We examine a tabular dataset $\mathcal{D} = \{(\mathbf{x}^i, \mathbf{y}^i)\}_{i=1}^{N}$ composed of $N$ annotated instances, where each input vector $\mathbf{x}^i$ encompasses $d$ features, and their respective names in natural language are given as $F = \{f_j\}_{j=1}^{d}$, with explicit definitions available. Within the scope of classification, each label $\mathbf{y}^i$ belongs to a specified class set. For $k$-shot learning scenarios, a subset of $k$ samples ($k < N$) is randomly selected from the dataset to serve as training data.

\subsection{Prompt Design for Extracting Rules}
\label{Sec:prompt_design}
To support the LLM in capturing rules through a robust reasoning methodology, we simulate the analytical strategies that expert practitioners employ for tabular data tasks. Our prompt is structured around three core elements (refer to Fig.~\ref{fig:prompt_for_rule} and Appendix~\ref{sec:example_rule}):

\vspace*{-0.12in}
\paragraph{Basic information description.} This section delivers foundational context necessary for the task (highlighted as orange text in Fig.~\ref{fig:prompt_for_rule}). It consists of: a task overview and associated label details, feature descriptions, and illustrative examples. The task description is typically framed as a query (for instance, ``Does this patient's myocardial infarction complications data indicate chronic heart failure? Yes or no?"), with additional examples in Appendix~\ref{sec:task_desc}. Feature descriptions convey the type of value and optionally the definition. For categorical features, category examples may be enumerated. Training instances are serialized as textual examples, including their actual labels, using formats adopted in prior works~\cite{dinh2022lift,hegselmann2023tabllm}:

\begin{align}
&\text{Serialize}(\mathbf{x}^i,\mathbf{y}^i, F) = \nonumber \\ 
&``\ f_1\ \text{is}\ \mathbf{x}^i_1.\ ... \ f_d\  \text{is}\  \mathbf{x}^i_d.\ \ \text{Answer:}\  \mathbf{y}^i\ ”
\end{align}

\vspace*{-0.12in}
\paragraph{Reasoning instruction.} Our strategy seeks to improve the LLM's reasoning abilities by incorporating explicit reasoning instructions (see blue annotations in Fig.~\ref{fig:prompt_for_rule}). The reasoning instructions start with an opening sentence that aligns with the chain-of-thought paradigm~\cite{wei2022chain}. Subsequently, we organize the LLM's reasoning workflow into two distinct stages. In the initial stage, the LLM is tasked with assessing the causal relationships or behavioral tendencies between features and the task objective, deliberately avoiding reference to example demonstrations and instead drawing from its broader general knowledge or commonsense reasoning. This enables the LLM to maximally utilize its existing conceptual understanding of the problem. In the following stage, the LLM incorporates the information gleaned from example demonstrations together with insights from the first step to construct class-specific rules. By segmenting reasoning in this manner, the process discourages the model from leveraging misleading correlations found in irrelevant columns and helps sharpen its focus on relevant and important features. To regulate the complexity and scope of extracted rules—thus guarding against both underfitting and overlarge prompts—we stipulate a fixed quantity (10) of rules to be generated for each class\footnote{Further discussion regarding rule number is provided in Section~\ref{sec:analysis}.}. In rule generation, we instruct the LLM to consult the feature type details presented in the basic information section, minimizing the risk of type mismatches in the resulting rules.

\vspace*{-0.12in}
\paragraph{Response instruction.} To streamline downstream parsing and application of the deduced rules, we offer targeted instructions that direct the LLM’s formatting choices (see yellow text in Fig.~\ref{fig:prompt_for_rule}). The prompts delineate the expected layout for responses specific to each answer class (i.e., response formatting) and the preferred structure for individual rules (i.e., rule formatting). By following these guidelines, the LLM is able to select suitable templates that correspond to each feature’s type. In the absence of explicit directions, the LLM retains the flexibility to combine several rules using logical connectors such as AND or OR. A sample output is illustrated in Appendix~\ref{sec:outcome-rule}.

\subsection{Inferring Class Likelihood via Rules}
\label{Sec:training}

\paragraph{Feature creation via rule parsing.} The rules produced during the earlier phase serve as the basis for constructing fresh binary features. For each class, these features are assigned a value depending on whether a sample fulfills the relevant rule criteria for that class (i.e., values of '0' or '1'). Such features facilitate estimation of a sample’s membership probability for each class. Due to the fact that rules formulated by the LLM are expressed in natural language, an automated parsing mechanism must be established to systematically generate these features. Despite efforts to standardize the rule formatting during rule extraction to enhance parsability, the outputs generated by the LLM often exhibit inconsistent syntax and unanticipated textual variations~\cite{kovavc2023large}, such as replacing symbolic notation like ``$>$'' and ``$<$'' with verbal expressions ``greater than'' or ``less than,'' further complicating the parsing task.

To overcome these text parsing obstacles, we do not develop intricate code solutions, but instead employ the LLM for this step. The LLM excels at interpreting the intended semantics even in the presence of diverse syntactic formulations, and is also capable of generating straightforward code from instructional prompts, owing to its code-oriented training data. Accordingly, we prepare prompts containing the function names, descriptions of inputs and outputs, as well as the parsed rules, and provide these to the LLM. As illustrated in Figures~\ref{fig:prompt_for_parsing} and ~\ref{sec:example_parsing}-\ref{sec:example_parse_outcome} in the Appendix, these prompt templates and their resulting outputs are demonstrated. Once the LLM produces the relevant code, it is executed with Python's \texttt{exec()} method, enabling the conversion of the original data.

\vspace*{-0.12in}
\paragraph{Estimating class probabilities.} With new rule-derived binary features for each class, a straightforward way to evaluate the likelihood of a sample belonging to each class is by tallying the number of rules it satisfies (the sum of binary indicators for that class). Nonetheless, the impact or relevance of individual rules and features can vary, so their importance must be inferred from training data. To achieve this, we apply a standard machine learning (ML) algorithm to the binary features representing each class, allowing for the learning of feature weights. For instance, a bias-free linear model can be utilized. Let $\mathbf{z}_k^i \in \{0, 1\}^R$ denote the binary feature vector produced for class $k$ from sample $\mathbf{x}^i$, where $R$ signifies the rule count per class and $c$ the total number of classes. The corresponding class probability $\mathbf{p}^i$ is calculated as follows:

\begin{align}
\text{logit}_k^i &= \max(\mathbf{w}_k, 0) \cdot \mathbf{z}_k^i, \nonumber \\
\mathbf{p}^i &= \text{Softmax}({[\text{logit}_{1}^i, ..., \text{logit}_{c}^i]}). \label{eq:infer_prob}
\end{align}

In this context, $\mathbf{w}_k$ denote the tunable parameters within the linear model, representing the contribution of each feature. By constraining these weights to reside in the positive domain, we guarantee that features derived by the LLM do not detract from a class's probability during the learning process. The resulting logit vector undergoes transformation to class probabilities via the softmax operator. The model parameters $\mathbf{w}_k$ are optimized by minimizing the cross-entropy loss on the limited available training examples.

\vspace*{-0.12in}
\paragraph{Ensembling with bagging.} As a final step, we run the described methodology multiple times to construct a set of inference models whose predictions are aggregated using ensembling. With trained parameters from $T$ distinct runs, denoted as $\{\mathbf{w}[t]\}_{t=1}^T$, the ultimate class prediction is determined by averaging the output probability vector $\mathbf{p}^i$ obtained for each weight $\mathbf{w}[t]$, as expressed in Eq.~\ref{eq:infer_prob}.

To foster diversity in the ensemble's rule generation, a relatively high temperature setting is applied during LLM inference. Additionally, the sequence of few-shot examples is randomly shuffled, motivated by evidence that example ordering can affect LLM output~\cite{lu2022fantastically}\footnote{Further analysis on rule variety appears in Appendix~\ref{sec:diversity}}. To better harness predictive diversity and enhance reliability, we employ bagging: in each trial, a random subset of features or instances is sampled, an approach increasingly beneficial as the dataset's size grows. The ensemble procedure confers two key benefits. First, if the LLM produces faulty rules in particular runs, correct rules from other trials can counterbalance the errors, thereby bolstering robustness. This advantage leverages the self-consistency property of LLMs~\cite{wang2022self}, since rules repeatedly discovered across several trials are statistically more dependable. Second, ensembling directly mitigates the prompt-length restrictions of LLMs. For cases where tabular inputs surpass the maximal allowed prompt, bagging narrows their size, preserving performance integrity. While any single rule set may miss relationships involving certain features, assembling an array of weak models across trials enables comprehensive coverage and neutralizes omissions present in individual runs.

\section{Experiments}
We assess the effectiveness of \textsf{FeatLLM} on a diverse collection of tabular datasets, performing a series of studies to better understand the mechanics of our framework. Owing to limitations in manuscript length, further experimental results—including LLM ablations, qualitative assessments, and additional analyses—are provided in the Appendix.

\subsection{Performance Evaluation}
\paragraph{Datasets.}
Our empirical studies cover 13 distinct datasets tailored for binary and multi-class classification tasks. The datasets include: (1) Adult~\cite{asuncion2007uci}, tasked with predicting if an individual’s annual income exceeds \$50,000; (2) Bank~\cite{moro2014data}, aimed at forecasting whether a customer subscribes to a term deposit; (3) Blood~\cite{yeh2009knowledge}, for identifying likelihood of repeat donations from blood donors; (4) Car~\cite{kadra2021well}, used to rate car quality; (5) Communities~\cite{misc_communities_and_crime_183}, for estimating local crime rates; (6) Credit-g~\cite{kadra2021well}, where the objective is to predict an individual’s creditworthiness; (7) Diabetes\footnote{\texttt{kaggle.com/datasets/uciml/pima-indians-diabetes-database}}, focused on classifying the presence or absence of diabetes; (8) Heart\footnote{\texttt{kaggle.com/datasets/fedesoriano/heart-failure-prediction}}, centering on the likelihood of coronary artery disease; and (9) Myocardial~\cite{misc_myocardial_infarction_complications_579}, which labels chronic heart failure patients.

Additionally, our benchmarking suite encompasses recent additions and artificially constructed datasets seldom addressed in prior literature or LLM pre-training: (10) Cultivars\footnote{\texttt{archive.ics.uci.edu/dataset/913}}, which classifies grain yield categories of soybean cultivars; (11) NHANES\footnote{\texttt{archive.ics.uci.edu/dataset/887}}, set up for age group prediction from medical records; (12) Sequence-type, which differentiates between arithmetic, geometric, Fibonacci, and Collatz sequences; and (13) Solution-mix, aimed at determining whether the solution mixture’s concentration surpasses 0.5. Cultivars and NHANES were submitted to the UCI Machine Learning Repository following September 2023, ensuring their absence from the LLM API’s pre-training corpus (specifically, gpt-3.5-turbo-0613). The final two are synthetic datasets to preclude LLM memorization effects and thus offer unbiased evaluation.

Table~\ref{Tab:basic-desc} summarizes key statistics demonstrating heterogeneity in dataset scale and intricacy. Every dataset supplies explicit feature names and accompanying descriptions. For instances such as `Communities' and `Myocardial', the attribute count exceeds 100, amplifying the challenge given LLMs’ finite input contexts.

\vspace*{-0.12in}
\paragraph{Baselines.}
We benchmark our method against nine established approaches. The initial group consists of classic algorithms for tabular data: (1) Logistic regression (LogReg), (2) XGBoost~\cite{chen2016xgboost}, and (3) RandomForest~\cite{ho1995random}. The subsequent group leverages unlabeled data: (4) SCARF~\cite{bahri2022scarf} and (5) STUNT~\cite{nam2023stunt}. Acquiring expansive unlabeled datasets can often be impractical in realistic deployment settings. A further competing technique is (6) TabPFN~\cite{hollmann2022tabpfn}, which relies on synthetic data of varying distributions to support model pre-training. The final cohort employs LLM-based strategies: (7) In-context learning (In-context)~\cite{wei2022emergent}, (8) TABLET~\cite{slack2023tablet}, and (9) TabLLM~\cite{hegselmann2023tabllm}. In-context learning attaches a handful of labeled instances directly in the prompt, abstaining from additional model adjustments. TABLET enriches the input by introducing rule-based and prototype-driven cues from an external classifier. TabLLM implements parameter-efficient adaptations such as IA3~\cite{liu2022few}, utilizing small data scenarios to update LLM behavior.

\vspace*{-0.12in}
\paragraph{Implementation details.}
\textsf{FeatLLM} is built upon GPT-3.5 as its foundational LLM, though the framework does not depend on this specific model (see Appendix~\ref{sec:backbones} for performance comparisons with alternate backbones). During inference, the temperature parameter for the LLM is established at 0.5, while top-p remains at the API's default setting of 1. For ensemble operations and rule extraction, we use 20 and 10 as their respective quantities. Insights regarding the influence of hyper-parameters are depicted in Figure~\ref{fig:hyperparameter2} and further discussed in Appendix~\ref{sec:hyper-parameter}.
A linear model is optimized via Adam with a learning rate of 0.01, trained across 200 epochs. The selection of the best epoch leverages $k$-fold cross-validation. When re-implementing baseline approaches, in-context learning utilizes GPT-3.5, whereas TabLLM adopts T0~\cite{sanh2022multitask}, consistent with its original specifications. Note that TabLLM exclusively permits LLMs for which checkpoints are openly released, making GPT-3.5 unavailable. For standard machine learning algorithms including LogReg, XGBoost, and RandomForest, we identified their hyper-parameters through a Grid-search and $k$-fold cross validation process. The parameter $k$ was set to either 2 or 4 to ensure every class appears in the training split. Additional implementation specifics can be found in Appendix~\ref{sec:baselines}.

\vspace*{-0.12in}
\paragraph{Results.}
Tables~\ref{tab:main1} and \ref{tab:main2} provide a comprehensive comparison of performance across multiple datasets. Each experiment is executed three times with varying random splits for the training and test sets, and we report the mean and standard deviation of the resulting AUC scores. The \textsf{FeatLLM} framework demonstrates leading or near-leading effectiveness relative to all baseline methods. Remarkably, our approach attains results comparable to those achieved by classical tabular baselines operating with 32 training samples, while only utilizing 4-shot examples (refer to Fig.~\ref{fig:compares}). Although models like STUNT and TabPFN exhibit substantial improvements on selected datasets, their aggregate gains over traditional machine learning algorithms remain limited. Furthermore, LLM-driven methods, including in-context learning, TABLET, and TabLLM, are unable to perform inference effectively when faced with tabular data of high dimensionality. Conversely, our method, which integrates bagging and ensemble strategies, consistently exhibits robust results even on high-feature-count datasets. It is worth mentioning that the bagging and ensemble scheme employed can theoretically be transferred to other LLM-based methods. However, this would necessitate conducting twice as many inferences per data point, substantially increasing computational requirements and thereby hindering practical deployment.

\subsection{Analysis \& Discussion}
\label{sec:analysis}
\paragraph{Ablation study.} We investigate the individual impact of primary framework components through a series of ablation studies, where the following variants are considered: (1) \textsf{-Tuning}: eliminates the linear model's weight adjustment, instead aggregating binary features directly to compute logits; (2) \textsf{-Ensemble}: forgoes the use of ensembling; (3) \textsf{-Description}: dispenses with incorporating feature descriptions; (4) \textsf{-Reasoning}: omits the initial reasoning instruction step during rule generation.

The outcomes, summarized in Table~\ref{tab:ablation}, reflect the average reduction in AUC across all datasets for each ablated variant. Performance consistently deteriorates whenever a core component is removed. Notably, the advantage conferred by weight tuning intensifies as the shot count rises, implying that richer data contexts permit more precise rule weighting and underscore the value of tuning. Conversely, the effects of incorporating feature descriptions and reasoning instructions are most prominent in low-shot scenarios, pointing to the crucial role that leveraging LLM’s prior knowledge plays when data is sparse. Finally, ensemble methods uniformly enhance results, underscoring their importance across all experimental setups.

\vspace*{-0.12in}
\paragraph{Training \& inference time comparison.} To evaluate computational efficiency, we assess both training and inference times across various methods. Our experiments utilize the Adult dataset with a training set comprising 16 shots. All methods are run on a single A100 GPU unless noted otherwise; TabLLM employs four A100 GPUs to enable model parallelism. Table~\ref{tab:cost} lists the comparative runtimes. Importantly, \textsf{FeatLLM} achieves an inference time on par with traditional baselines, maintaining low latency. Conversely, other LLM-based strategies, including In-context learning, TABLET, and TabLLM, incur substantially greater inference time. When examining training time, \textsf{FeatLLM} aligns closely with competing LLM-based techniques: it utilizes just 30 API calls, which are lightweight in terms of cost and can be executed in parallel, thereby mitigating resource constraints.

\vspace*{-0.12in}
\paragraph{Quality of features generated by \textsf{FeatLLM}.}
We carry out a straightforward experiment aimed at assessing how informative the rule-based features produced by \textsf{FeatLLM} are for real-world tasks. In this assessment, we calculate the mean AUROC between each generated feature and the target labels. Subsequently, we determine what proportion of features have an AUROC exceeding 0.5 for each dataset. 
The summarized outcomes are presented in Table~\ref{tab:quality}. The findings reveal that a significant portion of the rules devised are pertinent to the target variable, offering valuable information for task resolution (see the ``Before tuning'' column). Additionally, as a linear model is fitted to the data, rules with marginal utility (defined as those with learned weights beneath a specified threshold) are systematically excluded (refer to the ``After tuning'' column). The weight threshold is set to 0.1, in accordance with the overall weight histogram.

\vspace*{-0.12in}
\paragraph{Impact of introducing spurious correlations.} To evaluate how well different approaches handle noisy and irrelevant correlations under limited labeled data, we performed a targeted experiment. Specifically, we used the Heart dataset and supplemented it with randomly chosen columns from the Adult dataset, which bear no relevance to the Heart prediction task. Model performance was then tracked as we incrementally added these extraneous columns. To ensure fairness, all methods were provided with exactly 8 labeled samples and no access to any additional unlabeled data, thereby omitting approaches such as SCARF and STUNT from the comparison.

Figure~\ref{fig:spurious} displays the effects that spurious correlations have on model performance. Notably, \textsf{FeatLLM} exhibited the smallest drop in accuracy among all tested baselines when exposed to irrelevant features. This resilience is likely attributable to the system's explicit reliance on prior knowledge to identify meaningful relationships between the target task and candidate features during rule extraction; as such, rules built upon extraneous columns are avoided, enhancing robustness. We found that rules involving noisy columns made up only 1-2\% of the generated rules. However, \textsf{FeatLLM} is not immune to pitfalls—when columns originating from semantically similar datasets are appended, the model may still inadvertently learn spurious associations and misinterpret causal links, as discussed in Appendix~\ref{sec:spurious-appendix}.

\vspace*{-0.12in}
\paragraph{Hyper-parameter analysis}
\textsf{FeatLLM} utilizes two primary hyper-parameters: the number of ensembles incorporated and the quantity of rules per class generated with each LLM call. To evaluate their effects on model effectiveness, we systematically varied these hyper-parameters in our experiments. We observed that augmenting the ensemble count generally enhances predictive accuracy; however, the incremental gains tend to plateau once the ensemble size surpasses approximately 20 (refer to Fig.~\ref{fig:appendix-hyperparameter1} in the Appendix). Regarding the number of extracted rules, the experimental variations did not lead to any statistically significant differences. Nevertheless, we found that selecting 10 rules offers a favorable balance between prompt complexity and predictive efficacy (see Fig.~\ref{fig:hyperparameter2}). We hypothesize that expanding the rule count increases the dimensionality of the input, which can introduce additional information but may also promote overfitting, particularly within the few-shot learning scenario.

\section{Conclusion}
This work introduces \textsf{FeatLLM}, a novel framework that elevates performance in LLM-driven few-shot tabular data analysis. Distinct from methods that solely employ LLMs for instance-level prediction, \textsf{FeatLLM} targets the derivation of predictive heuristics, mirroring the strategies of feature engineering. By leveraging rule induction together with an ensemble-based bagging methodology, \textsf{FeatLLM} enables efficient inference, requires only API-level model interaction (with no fine-tuning necessary), and circumvents restrictions related to data dimensionality. The framework consistently demonstrates robust accuracy and offers a compelling alternative for practitioners seeking data-efficient solutions for real-world tabular applications utilizing LLMs.

\textsf{FeatLLM} has thus far been oriented towards low-shot learning contexts. Looking ahead, our research agenda involves extending its feature generation capacities for larger datasets, experimenting with alternative forms of features beyond simple rule extraction, and enhancing interpretability to facilitate broad utility across a spectrum of machine learning tasks.

\section{Broader Impact} The introduced \textsf{FeatLLM} framework leverages LLMs’ expansive prior knowledge in tabular learning tasks, thereby advancing beyond traditional supervised methods in terms of achievable accuracy. By increasing data efficiency, our approach offers mechanisms to address overfitting challenges associated with small labeled datasets, complementing limited data with the LLM's background knowledge. This is particularly advantageous for practitioners in fields such as finance or healthcare, where annotating data is expensive or difficult, and LLMs’ pretrained expertise in these domains can provide significant added value. However, as \textsf{FeatLLM} draws heavily on model prior knowledge that may influence and potentially override patterns in the collected labels, thorough assessment of social biases and misinformation embedded within these priors will be essential. Addressing these challenges remains an important direction for future research toward the responsible deployment and adoption of such systems.

\end{document}